\subsection{局部性：测量与聚合的顺序}
表~\ref{tab:component_ablation}按三个问题组织匹配控制。
块A显示，Local单独为0.83/0.84，Global为0.86/0.87，
二者融合达到0.88/0.88。
局部分支的作用体现为对整体证据的补充，而不是单独替代Global。

块B检验这一补充是否需要保留逐位置测量。
先平均每帧Patch、再计算D2的pooled控制为0.85/0.85。
为排除拟合向量数量的影响，我们将D2的拟合位置数匹配到pooled的预算：
每片段保留$14K_{\rm fit}$个抽样位置，至多42个，测试观察不变。
结果为0.88/0.88，相对pooled提高0.02/0.03，
差值区间分别为$[0.015,0.025]$与$[0.015,0.026]$。
CDF前的Local原始分数也提高0.03/0.03，两项区间均为正。
这与式~\eqref{eq:dispersion}保留位置离散信息的动机一致；
ComGenVid上pooled更好，则表明局部粒度的收益仍取决于数据域。

\subsection{二阶方向与标量替代}
块C固定目标Global、观察窗口和真实参考身份，每个候选重建自身统计与CDF。
D2相对D1提高Average AUC/AP 0.01/0.01，两项配对区间均为正；
23个单元中，AUC有20个、AP有22个取得正增量。
该结果支持当前真实评分框架中的阶次选择。

匹配Patch输入下，TTR、SPLIT统计及二维(TTR, LSMI)Gaussian控制的Average低于方向模型。
二维控制也使用相同目标真实片段，进一步区分增加真实统计拟合与保留高维方向的作用。
这些是同Global下的替代统计，不是SPLIT原生编码器的完整复现。
优势并非逐生成器一致：D2相对TTR融合的AUC/AP有12/11个单元胜出；
VideoFeedback的ModelScope单元中，TTR为0.96/0.97，高于D2的0.80/0.80。

\subsection{真实参考的重复验证}
为检验增量是否依赖一次真实抽样，我们从相同真实来源另选2357个新真实视频，
与原拟合、评价及百分位参考的已知源组互斥。
每域建立五个200片段拟合池，同次Global与Local共享视频并分别重建CDF。
五池之间允许重叠；表~\ref{tab:reference_stability}报告逐次计算指标后的均值。

新池中Full相对Global的Average AUC提高0.01，区间为$[0.005,0.015]$；
AP提高0.01，但区间$[-0.001,0.009]$跨零。
D2相对D1的两项增量为0.01/0.01，区间均为正。
因此，对Global的AUC补充与二阶相对一阶的收益得到重复支持，
而AP增量仍受参考条件影响。

\subsection{观察与重编码}
将FC K3换为Uniform K3，完整模型的Average由0.88/0.88变为0.86/0.87。
Uniform下加入Local的增量区间跨零，表明局部收益与观察位置有关。
48个16帧观察视频的计时子集中，FC K3在ComGenVid、VideoFeedback和GenVideo分别为
0.60、0.19和0.38秒，Uniform K3为0.53、0.14和0.31秒；
这些耗时包含解码、粗扫及评分，动态观察并非无成本提升。

在预定1534个片段身份上，固定原FC窗口和干净参考，对真假视频同样施加H.264重编码
（CRF23/35，4:4:4，不缩放或补帧）。
完整模型在原始、CRF23与CRF35条件下分别为
0.88/0.88、0.88/0.88和0.86/0.86。
较重编码使AUC/AP下降0.01/0.02，两项区间均为负。
Local减轻了该设置下的下降幅度，但压缩后的直接Full--Global增量区间跨零，
因此不将其解释为普遍的压缩鲁棒性。
